\documentclass[11pt]{article}

\usepackage[margin=1in]{geometry}
\usepackage{amsmath,amssymb}
\usepackage{graphicx}
\usepackage{booktabs}
\usepackage{hyperref}
\usepackage{authblk}
\usepackage[T1]{fontenc}
\usepackage[utf8]{inputenc}
\usepackage{natbib}
\usepackage{caption}

\hypersetup{
  colorlinks=true,
  linkcolor=blue,
  urlcolor=blue,
  citecolor=blue
}

\title{Task-Specific Computational Advantage of the \emph{Drosophila}
Connectome: A Null-Model-Controlled Study}

\author[1]{Abdulkadir Özcan}
\affil[1]{Independent Researcher \\ \texttt{a.kadir.ozcan@icloud.com}}

\date{September 2026}

\begin{document}
\maketitle

\begin{abstract}
Does the wiring topology of a real biological neural network -- who is
connected to whom, and with what synaptic weight -- confer a
computational advantage over a randomly rewired network with the exact
same size and degree distribution? We test this question using the
FlyWire FAFB v783 connectome of the adult female \emph{Drosophila
melanogaster} brain, embedded as a fixed recurrent substrate in a
differentiable, Dale's-law-preserving leaky-rate neural network
(\texttt{RateBrain}), trained by gradient descent on a variety of
computational tasks and compared against a family of null models that
progressively relax structural constraints (degree sequence, modularity,
weight identity). We find that the real connectome provides a
statistically robust advantage (Cliff's $\delta = 0.61$--$1.00$, $n=30$,
$p < 0.05$ after Bonferroni-style multiple-testing discipline via a
pre-registered magnitude-triggered extension rule) on a specific task
format: single-shot, feed-forward, continuous spatial/physical direction
regression. This holds across six independent physical/engineering
design problems spanning distinct failure modes (geotechnical slip,
bending stress, transverse shear, elastic buckling, deflection
serviceability, and thin-wall hoop stress), across an order-of-magnitude
range of subgraph sizes, and across a $50\times$ range of input sensing
scale. The advantage is largely (but not fully) explained by the
network's modular organisation, does not generalise to abstract
optimisation, discrete classification, closed-loop control, a different
species (\emph{C.~elegans}), a different individual of the same species
(male \emph{Drosophila} CNS), input noise, or weight quantisation, and
does not transfer zero-shot to a mechanically different task in the same
family. We report the full boundary-mapping process transparently,
including two methodological errors we made and corrected during the
study (a ratio-metric artifact and two term-imbalance calibration bugs
in newly designed benchmarks), as a case study in the discipline required
for this kind of null-model-controlled connectome research.
\end{abstract}

\section{Introduction}

A recurring question in computational neuroscience and connectome-based
reservoir computing \citep{suarez2024conn2res, costi2025flywire} is
whether the specific topology of a mapped biological neural circuit
does any real computational work beyond what is explained by simpler
statistical properties of the network (size, density, degree
distribution). Most prior work in this space reports a positive result
against a single, often permissive, null model and does not
systematically probe the boundaries of the finding. Our contribution is
methodological as much as empirical: we test one connectome, with one
consistent statistical framework, across a deliberately wide spectrum of
task types, and report negative and positive results with equal rigor,
including corrections to our own errors made mid-study.

Our central hypothesis (task--architecture matching) is that any
advantage of the real connectome over a randomised network should be
largest for tasks structurally close to what the connectome evolved to
do -- filtering a large, dense sensory input down to a small number of
motor outputs -- and should weaken or vanish for tasks further from
that structural niche.

\section{Methods}

\subsection{Data}

The primary connectome is FlyWire FAFB v783 \citep{dorkenwald2024flywire},
the electron-microscopy-reconstructed connectome of an adult female
\emph{Drosophila melanogaster} brain (139{,}255 neurons, 15{,}091{,}983
directed synaptic connections), with cell-type annotations from
\citet{schlegel2024flywire}. Comparison connectomes are the
\emph{C.~elegans} hermaphrodite somatic nervous system
\citep{cook2019celegans} and a male \emph{Drosophila} CNS connectome
obtained from an independent third-party dataset.

\subsection{Subgraph construction}

From the full connectome, we select encode (afferent-pool) and decode
(efferent-pool) neurons verified, via multi-hop breadth-first search on
the full directed graph, to have a genuine forward path between them
(a naive random afferent/efferent draw can have zero true connectivity).
From this seed set we grow an induced subgraph of a target size (the
official size is 3{,}000 nodes; robustness across 1{,}000--10{,}000 nodes
is reported in Section~\ref{sec:generality}) via breadth-first search
along real directed edges.

\subsection{Substrate and training}

\texttt{RateBrain} is a differentiable, leaky-rate recurrent network
with Dale's-law-preserving weights ($w = \mathrm{sign} \cdot
\mathrm{softplus}(\mathrm{gain})$), trained by standard backpropagation
(Adam optimiser). For single-shot tasks, sensory input (finite-difference
``rays'' -- local probes of a task's objective function along each
design-variable axis) is injected into encode neurons and held constant
for $T{=}8$ recurrent sub-steps; a linear readout from decode neurons
produces a predicted direction. The teacher signal is a crude
finite-difference negative-gradient estimate from the same probes (not a
hand-crafted heuristic), and training minimises mean squared error
against it.

\subsection{Null model family}

We compare the real connectome against six null models of increasing
structural fidelity:
\begin{itemize}
  \item \texttt{er\_null} -- node/edge count only;
  \item \texttt{scale\_free\_null} -- generic heterogeneous degree distribution;
  \item \texttt{scale\_free\_correlated\_null} -- plus in/out-degree correlation;
  \item \texttt{weight\_shuffle} -- full topology, shuffled weights;
  \item \texttt{community\_preserving\_rewire} -- full degree sequence plus full modularity/community structure;
  \item \texttt{degree\_preserving\_rewire} -- exact degree sequence via Maslov--Sneppen double-edge-swap rewiring (our primary/official null model throughout).
\end{itemize}

\subsection{Statistics}

We use the triad of Wilcoxon signed-rank (paired), Mann--Whitney~$U$
(unpaired), and Cliff's~$\delta$ (effect size) throughout. The official
gate is $p<0.05$ \emph{and} $|\delta|>0.33$. The standard protocol is an
$n=8$ pilot; whenever $|\delta|>0.33$ at $n=8$, the test is extended to
$n=30$, \emph{independent of whether the pilot $p$-value was already
significant} -- a magnitude-triggered, not significance-triggered,
extension rule, decided in advance and applied uniformly.

\section{Results}

\subsection{The official finding}

Two independent physical/engineering tasks -- slope-stability
(geotechnical limit-equilibrium factor of safety) and cantilever beam
design (bending stress) -- show a strong, reproducible advantage of the
real connectome over the \texttt{degree\_preserving\_rewire} null at
$n=30$: $\delta = 0.884$ ($p < 10^{-4}$) and $\delta = 0.613$
($p<10^{-4}$) respectively. The advantage replicates in 9 of 11
independent subgraph seeds (median $\delta = 0.500$), so it is not an
artefact of one fortunate subgraph draw.

\subsection{Mechanism}

Sweeping the null-model family for the slope-stability task
(Table~\ref{tab:mechanism}) shows that $\delta$ falls from $0.884$
(degree-preserving) to $0.429$ once the null is also constrained to
match the real network's community/modularity structure
(\texttt{community\_preserving\_rewire}), and becomes non-significant
once weights alone are shuffled on the real topology
(\texttt{weight\_shuffle}, $\delta = 0.258$, $p=0.088$). Modularity thus
explains roughly half of the effect; the remainder is unexplained
(a weak, uncorrected exploratory correlation with the fraction of
\emph{ascending}-type interneurons was noted but not confirmed).

\begin{table}[h]
\centering
\caption{Cliff's $\delta$ across the null-model family (slope stability,
$n=30$).}
\label{tab:mechanism}
\begin{tabular}{lccc}
\toprule
Null model & Preserves & $\delta$ & $p$ \\
\midrule
\texttt{degree\_preserving} & full degree sequence & $0.884$ & $\approx 0$ \\
\texttt{scale\_free} & generic heterogeneity & $0.738$ & $\approx 0$ \\
\texttt{scale\_free\_correlated} & + degree correlation & $0.604$ & $0.00006$ \\
\texttt{community\_preserving} & + full modularity & $0.429$ & $0.0044$ \\
\texttt{weight\_shuffle} & topology, no weights & $0.258$ & $0.088$ (n.s.) \\
\bottomrule
\end{tabular}
\end{table}

\subsection{Generality across the physical-task family}
\label{sec:generality}

We extended the official pair of tasks to a family of six mechanically
distinct single-shot physical/engineering design problems, each with its
own failure/design criterion (Table~\ref{tab:physical}). All six show a
statistically robust advantage at $n=30$, with zero exceptions.

\begin{table}[h]
\centering
\caption{Physical/engineering task family, all against
\texttt{degree\_preserving\_rewire}.}
\label{tab:physical}
\begin{tabular}{llcc}
\toprule
Task & Failure/design criterion & $\delta$ ($n{=}30$) & $p$ \\
\midrule
Beam deflection & serviceability (stiffness) & $1.000$ & $3\times10^{-11}$ \\
Slope stability & geotechnical limit-equilibrium & $0.884$ & $\approx 0$ \\
Pressure vessel & thin-wall hoop stress & $0.871$ & $7\times10^{-9}$ \\
Column buckling & elastic instability & $0.88$ & $5\times10^{-9}$ \\
Beam shear & transverse shear stress & $0.811$ & $7\times10^{-8}$ \\
Beam bending & bending stress & $0.613$ & $\approx 0$ \\
\bottomrule
\end{tabular}
\end{table}

Two of the six tasks (column buckling and pressure vessel) initially
returned an apparent \texttt{FAIL} and were provisionally reported as
such before we discovered, via a suspicious numerical similarity between
two independently designed benchmarks' diagnostic output, that both
suffered from an objective-term-balance calibration bug: in one case the
material-cost term dominated $>99.99\%$ of the objective at typical
sampled points (making the task trivially solvable by any network
regardless of structure), and in the mirror-image case a physical stress
term dominated $>97\%$. Both were corrected (narrowing the design-variable
box and/or re-deriving load magnitude as a fraction of a reference
cross-section's capacity, verified numerically before any further GPU
time was spent) and both then showed some of the strongest positive
results in the entire study. This episode is documented in full in the
project's public experiment log as a cautionary methodological case study
(Section~\ref{sec:limits}).

We also swept the size of the 3{,}000-node official subgraph across
1{,}000--10{,}000 nodes (a 10$\times$ range): all four sizes tested pass
the gate ($\delta = 1.00, 1.00, 0.84, 0.81$ respectively), so the official
finding is not an artefact of the specific 3{,}000-node choice. Separately,
zero-shot evaluation of the officially trained network at input sensing
radii $0.1\times$ and $5\times$ the training radius (a $50\times$ total
range) both pass the gate ($\delta = 0.684$ and $0.906$ respectively,
$n=30$/$n=8$), and mechanism analysis shows this generalisation axis
collapses under the exact same null model (\texttt{community\_preserving})
as the main finding -- i.e. it is not an independent second mechanism.

\subsection{Boundary: where the advantage does not hold}

The advantage does not generalise to: abstract, non-spatial optimisation
landscapes (Sphere, Rastrigin, Ackley, Rosenbrock, Beale; $\delta \approx
0$ to $-0.34$); a discrete-classification reformulation of the same
underlying task (octant classification, $\delta=0.141$); closed-loop,
multi-step control variants of the same task ($\delta$ falls to
$\approx 0$); a different species (\emph{C.~elegans}, all null models
n.s.); a different individual of the same species (male \emph{Drosophila}
CNS, $\delta=0.250$, n.s., confirmed by two independent follow-up
controls); or an external gradient-free search algorithm warm-started
with the network's predictions (PSO warm-start hybrid, $\delta=0.125$,
n.s.).

Two further axes show a \emph{reversed} advantage (real connectome
performing significantly \emph{worse} than the null): zero-shot transfer
of a slope-stability-trained network to the mechanically related
beam-design task without retraining ($\delta=-0.464$, $n=30$), and
robustness to moderate ($15\%$) proportional input noise ($\delta=-0.911$,
$n=30$, confirmed independently via both ratio- and absolute-error
metrics). By contrast, an initially reported reversed-advantage finding
for random structural damage (edge ablation) was retracted after
independent audit revealed it to be a ratio-metric artefact (the real
network's already-lower baseline error mechanically inflated a
ratio-based fragility metric); recomputed on absolute error, the
official 20\% ablation level shows no significant difference
($\delta=-0.176$, $p=0.246$).

Finally, the advantage evaporates under weight quantisation at every
tested precision (8-, 4-, and 2-bit, all non-significant, $n=8$--$30$),
indicating that it depends on full floating-point numerical precision
and should not be expected to transfer directly to low-bit neuromorphic
hardware.

\section{Discussion}

Taken together, the results support a narrow but well-evidenced claim:
a specific biological connectome can carry a structural computational
advantage for the class of computation it plausibly evolved for (dense
sensory input funnelled to a small motor output, computed in a single
feed-forward pass), and this advantage is (i) acquired from remarkably
little training data, (ii) robust to substantial changes in input
sensing scale and subgraph size, (iii) largely but not fully attributable
to modular network organisation, and (iv) fragile outside this narrow
task format -- to a different task class, a different species or
individual, input noise, and reduced numerical precision.

This is a materially narrower and more precisely bounded claim than
``biological networks are generally better computers.'' It is, however,
a substantially \emph{broader} claim within its own boundary than we
initially believed: what first appeared to be an idiosyncratic result on
two specific engineering problems turned out, once two calibration bugs
were found and fixed, to hold across six mechanically distinct physical
design criteria with no exceptions.

\subsection{Limitations and lessons}
\label{sec:limits}

All official findings derive from a single female specimen's connectome;
subgraph-seed sampling ($n=11$) supports but does not prove full
generality. The \texttt{RateBrain} architecture is a biologically
plausible but simplified abstraction (no spike timing, no short-term
plasticity). Five additional task categories attempted in an
overnight autonomous experiment batch (temporal integration, working
memory, anomaly detection, chaotic time-series prediction, time-to-contact
estimation) were found, on independent audit, to be invalid due to a
shared infrastructure bug (a mismatch between full-graph connectivity
verification and subgraph-extraction breadth-first search that can leave
a single encode neuron disconnected from all decode neurons within the
extracted subgraph) and yield no usable evidence either way; the root
cause is documented but the shared code was deliberately left unmodified
pending careful, isolated validation against the existing official
results.

We highlight, as a methodological contribution in its own right, that
two of the six physical-task results in Section~\ref{sec:generality}
were initially wrong in a diagnosable, correctable way, and that both
corrections were found only because an unrelated newly designed
benchmark's pre-training diagnostic numbers were suspiciously similar to
an already-reported result's diagnostics. We now treat a pre-training
check of (a) the fraction of the objective contributed by each additive
term at sampled points, and (b) the cosine similarity between the
finite-difference teacher signal and a near-exact local gradient
reference, as mandatory before any new single-shot task is trusted --
several of the exciting-looking early results in this line of work
(including, at different points, results in both directions) turned out
on inspection to be measurement artefacts rather than genuine findings
about the connectome, and only systematic self-auditing caught them.

\section{Conclusion}

The female \emph{Drosophila} FlyWire connectome's sensory-to-motor
subgraph shows a robust, mechanism-partially-explained, reproducible
computational advantage over degree-matched random networks, specifically
on single-shot, feed-forward, continuous spatial/physical direction
regression -- and this advantage now extends, following correction of
two calibration errors, across every physical engineering design
criterion we have tested. The advantage does not indicate general-purpose
computational superiority: it does not transfer across task format,
species, individuals, input noise, or numerical precision. All code, raw
results, and a complete, dated experiment log -- including every error
found and corrected during the study -- are publicly available at
\url{https://github.com/AutoPyloter/fly_op}.

\bibliographystyle{plainnat}
\begin{thebibliography}{9}

\bibitem[Suárez et al.(2024)]{suarez2024conn2res}
Suárez, L.E. et al. (2024).
Conn2res: A toolbox for connectome-based reservoir computing.
\emph{Nature Communications}.

\bibitem[Costi et al.(2025)]{costi2025flywire}
Costi, S., Hadjiivanov, A., Dold, D., Hale, J., Izzo, D. (2025).
The Drosophila connectome as a computational reservoir for time-series
prediction. \emph{Biomimetics}.

\bibitem[Dorkenwald et al.(2024)]{dorkenwald2024flywire}
Dorkenwald, S. et al. (2024).
Neuronal wiring diagram of an adult brain.
FlyWire FAFB v783, Zenodo DOI 10.5281/zenodo.10676866.

\bibitem[Schlegel et al.(2024)]{schlegel2024flywire}
Schlegel, P. et al. (2024).
Whole-brain annotation and multi-connectome cell typing of Drosophila.
\emph{Nature}.

\bibitem[Cook et al.(2019)]{cook2019celegans}
Cook, S.J. et al. (2019).
Whole-animal connectomes of both Caenorhabditis elegans sexes.
\emph{Nature}.

\end{thebibliography}

\end{document}
